% !TeX root = ../main.tex
\section{Observer transformations, gauges, and admissible data}
\label{sec:technical_setting}

This appendix records the technical setting assumed throughout the main
derivation: the superposed-observer transformation rules and the
frame-indifference restrictions they impose on the constitutive potentials, the
gauge conditions for the electrostatic and transfer potentials, and the
admissible initial and boundary data for a boundary-value problem.

\subsection{Frame indifference, material symmetry, and angular momentum}

Let a superposed rigid change of spatial observer be described by a translation
\(\mbf c(t)\) and a proper-orthogonal tensor \(\mbf Q(t)\).  Scalars are
observer invariant, apart from the gauges of the electrostatic and
Bedford--Drumheller transfer potentials specified below.  The kinematic,
electrical, flux, and stress variables used below transform as
%
\begin{subequations}
	\label{eq:observer_change_rules}
	\begin{align}
		\mbf x^\star
		 & =
		\mbf c(t)+\mbf Q(t)\mbf x,
		\qquad
		\mbf Q^T\mbf Q=\mbf I,
		\qquad
		\det\mbf Q=1,
		\label{eq:observer_position_rule} \\
		\mbf v_\xi^\star
		 & =
		\phasedot{\mbf c}+\phasedot{\mbf Q}\mbf x+\mbf Q\mbf v_\xi,
		\qquad
		\left(
		\mbf v_\xi-\mbf v_\zeta
		\right)^\star
		=
		\mbf Q
		\left(
		\mbf v_\xi-\mbf v_\zeta
		\right),
		\label{eq:observer_velocity_rule} \\
		\mbf F_\xi^\star
		 & =
		\mbf Q\mbf F_\xi,
		\qquad
		\mbf F^\star=\mbf Q\mbf F,
		\qquad
		\bar{\mbf F}_s^\star=\mbf Q\bar{\mbf F}_s,
		\qquad
		\mbf A_s^\star=\mbf Q\mbf A_s\mbf Q^T,
		\qquad
		s\in\mc S,
		\label{eq:observer_deformation_rules} \\
		\mbf h^\star(\mbf x^\star,t)
		 & =
		\mbf h(\mbf x,t),
		\label{eq:observer_capillary_history_rule} \\
		\varphi^\star(\mbf x^\star,t)
		 & =
		\varphi(\mbf x,t)+c_\varphi(t),
		\qquad
		\mbf E^\star=\mbf Q\mbf E,
		\qquad
		\mbf d_\xi^\star=\mbf Q\mbf d_\xi,
		\label{eq:observer_electrical_rules} \\
		\left(
		\mbf j_{\xi,k}^\alpha
		\right)^\star
		 & =
		\mbf Q\mbf j_{\xi,k}^\alpha,
		\qquad
		\mbf q_\xi^\star=\mbf Q\mbf q_\xi,
		\qquad
		\left(
		\bs\sigma_\xi^{\prime+}
		\right)^\star
		=
		\mbf Q\bs\sigma_\xi^{\prime+}\mbf Q^T .
		\label{eq:observer_flux_stress_rules}
	\end{align}
\end{subequations}
%
Here \(k\in\{{\rm disp},{\rm diff}\}\).  The entries of \(\mbf h\) are
internal scalar variables, so the bold notation collects them but does not make
\(\mbf h\) a spatial vector.  Consequently, the capillary-history-gradient
force in \eqref{eq:MC_overall_momentum_nonlinear_biot} transforms as a spatial
vector.
The transfer potential is a kinematic potential rather than an objective
thermodynamic scalar.  Between inertial observers, take
\(\mbf Q(t)=\mbf Q_0\) and
\(\mbf c(t)=\mbf c_0+t\mbf u_{\rm obs}\), where
\(\mbf u_{\rm obs}\) is constant in the starred frame.  Its Galilean
transformation is
%
\begin{equation}
	\tau^\star(\mbf x^\star,t)
	=
	\tau(\mbf x,t)
	+
	\mbf u_{\rm obs}\cdot\mbf x^\star
	-
	\frac{1}{2}
	\mbf u_{\rm obs}\cdot\mbf u_{\rm obs}\,t
	+
	c_\tau,
	\label{eq:transfer_potential_galilean_rule}
\end{equation}
%
where \(c_\tau\) is a time-independent gauge constant.  Consequently,
%
\begin{equation}
	\nabla_{\mbf x^\star}\tau^\star-\mbf v_\xi^\star
	=
	\mbf Q_0
	\left(
	\nabla_{\mbf x}\tau-\mbf v_\xi
	\right),
	\qquad
	\frac{D_\xi^\star\tau^\star}{Dt}
	-
	\frac{1}{2}
	\mbf v_\xi^\star\cdot\mbf v_\xi^\star
	=
	\frac{D_\xi\tau}{Dt}
	-
	\frac{1}{2}
	\mbf v_\xi\cdot\mbf v_\xi .
	\label{eq:transfer_potential_galilean_invariants}
\end{equation}
%
These are the two combinations that enter the insertion-momentum and
component-transfer equations.  A time-dependent rotation produces an observer
velocity that cannot be written as the gradient of a scalar, so this
transfer-potential form is restricted to inertial observers.  Noninertial
balance laws must include the corresponding fictitious forces.
Spatial gradients of objective scalars transform by left multiplication with
\(\mbf Q\).  The velocity gradient transforms as
\(\mbf L_\xi^\star=\phasedot{\mbf Q}\mbf Q^T+
\mbf Q\mbf L_\xi\mbf Q^T\), while its symmetric part transforms
covariantly.  Constitutive objectivity concerns arbitrary time-dependent
\(\mbf Q\); the balance laws themselves are written for inertial observers,
with fictitious accelerations included in the body force if a noninertial frame
is used.

The elastic factors are assigned consistently with
\eqref{eq:observer_deformation_rules}, so that
\(\bar{\mbf F}_s^{e\star}=\mbf Q\bar{\mbf F}_s^e\) and
\(\mbf A_s^{e\star}=\mbf Q\mbf A_s^e\mbf Q^T\).  Frame indifference of the
stored potentials therefore requires
%
\begin{subequations}
	\label{eq:stored_potential_objectivity}
	\begin{align}
		\psi_s\left(
		\mbf Q\bar{\mbf F}_s^e,
		\mbf Q\mbf A_s^e\mbf Q^T,
		\bar\rho_s,\{\eta_s^\alpha\},\theta_{\mc S}
		\right)
		 & =
		\psi_s\left(
		\bar{\mbf F}_s^e,\mbf A_s^e,
		\bar\rho_s,\{\eta_s^\alpha\},\theta_{\mc S}
		\right),
		\qquad
		s\in\mc S,
		\label{eq:solid_energy_objectivity} \\
		\omega_\xi^+\left(
		\mbf Q\mbf E,
		\bar\rho_\xi,\{\eta_\xi^\alpha\},\theta_\xi
		\right)
		 & =
		\omega_\xi^+\left(
		\mbf E,
		\bar\rho_\xi,\{\eta_\xi^\alpha\},\theta_\xi
		\right).
		\label{eq:electric_enthalpy_objectivity}
	\end{align}
\end{subequations}
%
The plastic and stress-free factors are intermediate-configuration variables;
their constitutive laws are understood to be invariant under a change of
intermediate basis.  Likewise, every tensor mobility, permeability, and thermal
conductivity used below must transform by
\(\mbf Q(\cdot)\mbf Q^T\), every vector-valued closure must be equivariant, and
the scalar reaction and plastic mobilities must be invariant.  If such a
transport tensor depends only on the scalar state variables displayed in this
manuscript, material isotropy reduces it to a scalar multiple of \(\mbf I\);
anisotropic transport requires an objective structural or resolved-flow
argument.

Because \(\omega_\xi^+\) in
\eqref{eq:electric_enthalpy_objectivity} has only one spatial-vector argument,
its derived electric displacement is parallel to the electric field.  Hence
%
\begin{equation}
	\mbf d_\xi(\mbf Q\mbf E,\ldots)
	=
	\mbf Q\mbf d_\xi(\mbf E,\ldots),
	\qquad
	\mbf E\times\mbf d_\xi=\mbf0,
	\qquad
	\bs\sigma_\xi^+
	=
	(\bs\sigma_\xi^+)^T .
	\label{eq:objective_electrical_response_symmetry}
\end{equation}
%
Thus the present electrical state set describes an isotropic, possibly
density-, composition-, and temperature-dependent dielectric response.  A
permanent polarization or bedding-directed dielectric response would require an
additional objective structural argument and the associated stress and torque
terms; those extensions are outside the present model.

The continuum is nonpolar and has no body couples, constituent couple stresses,
or interphase interaction couples.  Local angular-momentum balance therefore
imposes
%
\begin{equation}
	\bs\sigma_\xi^{\prime+}
	=
	(\bs\sigma_\xi^{\prime+})^T,
	\qquad
	\xi\in\mc F\cup\mc S .
	\label{eq:phase_angular_momentum_restriction}
\end{equation}
%
Equation~\eqref{eq:phase_angular_momentum_restriction} is a constitutive
restriction on the solid energy derivatives as well as on the electrical
stress.  Mechanical anisotropy remains admissible through the solid elastic
arguments and the assigned material-symmetry group.  An isotropic mechanical
specialization requires invariance under every proper material rotation and
therefore representation by objective scalar invariants and isotropic tensor
functions; no such specialization is assumed unless stated locally.

\subsection{Electrical ensemble and gauge}

The electrical specialization is electroquasistatic:
\(\mbf E=-\nabla_{\mbf x}\varphi\), with no electromagnetic-wave dynamics.
The phase electric enthalpies take \(\mbf E\) as their electrical argument, and
\(\mbf d_\xi=-\partial\omega_\xi^+/\partial\mbf E\).  Free bulk charge
\(\varrho\) and prescribed outward surface charge \(\bar q\) are source-work
coefficients rather than additional stored energies.  Material-carried charge
is obtained by weighting each component mass by \(z_\xi^\alpha\).

The boundary is partitioned into \(\Gamma_\varphi\), where \(\varphi\) is
prescribed, and \(\Gamma_q\), where
\(\mbf d\cdot\mbf n=\bar q\).  This electric-displacement datum is distinct from
normal material-current data.  For a pure electric-displacement
boundary problem,
\(\Gamma_q=\partial\mc B\), solvability and its fixed-domain rate require
%
\begin{subequations}
	\label{eq:electrical_pure_neumann_compatibility}
	\begin{align}
		\int_{\mc B}\varrho\,\dv
		 & =
		\int_{\partial\mc B}\bar q\,{\rm d}a,
		\label{eq:electrical_pure_neumann_gauss_compatibility} \\
		\int_{\partial\mc B}
		\left[
			\frac{\partial\bar q}{\partial t}
			+
			\mbf i\cdot\mbf n
		\right]{\rm d}a
		 & = 0.
		\label{eq:electrical_pure_neumann_rate_compatibility}
	\end{align}
\end{subequations}
%
If a nonempty potential boundary is present, its normal displacement adjusts
instead of being prescribed by these global equalities.  If no potential value
is prescribed, the additive gauge of \(\varphi\) must also be fixed, for example
by prescribing its mean.  Reaction sums conserve charge, and transport keeps
the chemical driving force separate from the electric field; neither
result changes when a constant is added to \(\varphi\).  At nonuniform
temperature, differentiating an electrochemical potential divided by
temperature would introduce that arbitrary constant.  The neutral thermal and
electric forces are therefore kept separate.  The reservoir
specialization takes
\(\omega_\xi^+\) to have the objective single-vector dependence in
\eqref{eq:electric_enthalpy_objectivity}, thereby excluding permanent
polarization at zero field, directional dielectric structure, and independent
dielectric relaxation.

\subsection{Initial and boundary data}

The local theory does not prescribe a single boundary-value problem.  An
application must supply a compatible independent set of initial
phase/component masses, active phase fractions, fluid and solid temperatures,
capillary-history variables, solid kinematic and internal variables, and, when
inertia is retained, phase velocities.  These data must satisfy the
mass-fraction, volume-fraction, mass, and charge restrictions, together with the
constitutive bounds on the capillary-history state.
The local evolution law for \(\mbf h\) requires initial data but no boundary
data; a nonlocal regularization would add the corresponding boundary
conditions.  Here \(s_\star\in\mc S\) denotes the selected solid-phase
transfer reference.  The transfer potential requires one
initial datum
for the skeleton-trajectory evolution
\eqref{eq:MC_admissible_conversion_component}.  The specific generalized transfer-work
potentials \(L_\xi^N\) of the non-reference phases are then recovered from
\eqref{eq:MC_phase_reference_transfer_work_recovery}, rather than prescribing
additional initial data for \(\tau\) along every phase trajectory.  The
transfer potential's time-independent additive constant may be fixed, for
example, by
%
\begin{equation}
	\tau(\mbf x,0)=\tau_0(\mbf x),
	\qquad
	\int_{\mc B}\tau_0\,\dv=0 .
	\label{eq:transfer_potential_initial_gauge}
\end{equation}
%
The electrostatic potential has no independent evolution datum; it is obtained
at each time from Gauss' law, the current charge state, and its boundary data.

Mechanical data use the complementary motion/traction partition defined with
\eqref{eq:W_def}.  The solid and fluid natural phase tractions are
\eqref{eq:phase_natural_traction} and
\eqref{eq:fluid_natural_traction}, respectively.  When the phase virtual motions
coincide on the mixture boundary, their boundary powers give the overall
traction \eqref{eq:overall_natural_traction}.  For every independent component
\(\alpha\ne N\) in phase \(\xi\in\mc G\), transport data use complementary
normal-flux and component-potential portions.  The electric-potential or
electric-displacement datum is supplied independently on
\(\Gamma_\varphi\) or \(\Gamma_q\):
%
\begin{subequations}
	\label{eq:component_boundary_data}
	\begin{align}
		\partial\mc B
		 & =
		\overline{\Gamma_{m,\xi}^{\alpha}}
		\cup
		\overline{\Gamma_{\overline{\mc M}_\xi^\alpha}},
		\qquad
		\Gamma_{m,\xi}^{\alpha}
		\cap
		\Gamma_{\overline{\mc M}_\xi^\alpha}
		=\varnothing,
		\label{eq:component_boundary_partition} \\
		\left(
			\rho_\xi^\alpha\mbf v_\xi
			+
			\mbf j_{\xi,{\rm disp}}^\alpha
			+
			\mbf j_{\xi,{\rm diff}}^\alpha
		\right)\cdot\mbf n
		 & =
		\bar m_{\xi}^{\alpha,{\rm out}}
		\qquad
		\text{on }\Gamma_{m,\xi}^{\alpha},
		\label{eq:component_normal_flux_boundary_data} \\
		\frac{\hat\mu_\xi^\alpha-\hat\mu_\xi^N}{\theta_{\mc G}}
		 & =
		\overline{
			\left(
			\hat\mu_\xi^\alpha-\hat\mu_\xi^N
			\right)/\theta_{\mc G}
		}
		\qquad
		\text{on }\Gamma_{\overline{\mc M}_\xi^\alpha}.
		\label{eq:component_potential_boundary_data}
	\end{align}
\end{subequations}
%
Impermeability sets
\(\bar m_{\xi}^{\alpha,{\rm out}}=0\).  The reference-component flux is
reconstructed rather than prescribed independently, and the relative parts of
the boundary fluxes must satisfy the zero-sum conditions in
\eqref{eq:averaged_relative_flux_sum}.

Each thermal subsystem uses a complementary temperature/heat-flux partition,
with the sign convention that positive prescribed flux leaves \(\mc B\).
Here \(\mbf q_\xi\) is the gauge-invariant nonadvective material-energy flux in
\eqref{eq:MC_energy_balance}; its Fourier closure is
\eqref{eq:MC_heat_flux_law}.  Electrical work on component transport relative
to a phase is a gauge-invariant volume power in
\eqref{eq:MC_energy_balance}, not an additional potential-dependent boundary
flux:
%
\begin{subequations}
	\label{eq:thermal_boundary_data}
	\begin{align}
		\partial\mc B
		 & =
		\overline{\Gamma_{\theta,\mc G}}
		\cup
		\overline{\Gamma_{q,\mc G}},
		\qquad
		\Gamma_{\theta,\mc G}
		\cap
		\Gamma_{q,\mc G}
		=\varnothing,
		\label{eq:thermal_boundary_partition} \\
		\theta_{\mc G}
		 & =
		\bar\theta_{\mc G}
		\qquad
		\text{on }\Gamma_{\theta,\mc G},
		\label{eq:temperature_boundary_data} \\
		\sum_{\xi\in\mc G}
		\mbf q_\xi\cdot\mbf n
		 & =
		\bar q_{\mc G}^{\rm out}
		\qquad
		\text{on }\Gamma_{q,\mc G},
		\qquad
		\mc G\in\{\mc F,\mc S\}.
		\label{eq:heat_flux_boundary_data}
	\end{align}
\end{subequations}
%
The corresponding entropy flux also contains the
component-potential-weighted component-relative contribution stated in
\eqref{eq:MC_entropy_inequality}.  Its normal value is fixed by the component
flux and potential data in \eqref{eq:component_boundary_data}, and it vanishes
on an impermeable boundary.  It is therefore not an additional independent
thermal boundary datum.

Insulation sets \(\bar q_{\mc G}^{\rm out}=0\).  For a steady pure-heat-flux
problem, integrating \eqref{eq:MC_energy_balance} shows that the prescribed
outward heat flux must equal the net nonconductive volume production minus the
nonconductive steady storage and transport terms.  This is the thermal
pure-Neumann compatibility condition, with its sign fixed by the
\(-\nabla_{\mbf x}\cdot\mbf q_\xi\) term in that balance.  Similarly, a
quasi-static pure-traction problem must satisfy the integrated overall momentum
balance \eqref{eq:MC_overall_momentum_nonlinear_biot}; with inertia retained,
the same integral determines the rate of total momentum.

Electrostatic data use the potential/displacement partition and compatibility
conditions \eqref{eq:electrical_pure_neumann_compatibility}.  Normal component
mass fluxes determine the material-carried current
\(\mbf i\cdot\mbf n\).  Pure-flux component data must also satisfy the component
mass balances integrated over \(\mc B\), so that initial mass plus integrated
mechanism supply minus outward boundary flux equals current component mass.
